\metatitle{Projections of assignment polytopes}
\metadescription{Arbitrary finite relations as projections of permutation
matroid simplices, and unbounded integer coefficients in three retained columns.}

\section[columnProjections]{Retaining color columns}

Retaining $r$ columns of an assignment polytope is different from working
with $r$ actual colors. The other colors can encode arbitrary relations
between the retained columns. We prove a universality theorem and then
give three-column projections with unbounded necessary integer coefficients.

Call a description \defin{binary} if its equations and upper
inequalities have coefficients in $\{0,1\}$. A \defin{ternary}
description allows coefficients in $\{-1,0,1\}$.

\section[projectionUniversality]{Arbitrary finite relations}

Let $I_1,\dotsc,I_r$ be disjoint finite sets and let
$\mathcal R\subseteq I_1\times\dotsb\times I_r$ be nonempty. With
standard basis vectors indexed by the disjoint union, write
\[
 P_{\mathcal R}=
 \operatorname{conv}\left\{\sum_{j=1}^r e_{i_j}:
 (i_1,\dotsc,i_r)\in\mathcal R\right\}.
\]

\begin{theorem}[Universality of column projections]
There is a permutation matroid whose projection onto $r$ selected
color columns is $P_{\mathcal R}$, after removing fixed zero coordinates.
Its full assignment polytope is a simplex with a complete binary
description in its assignment coordinates.
\end{theorem}

\begin{proof}
Choose a prime $p>\max(|\mathcal R|,r)$. On $\mathbb F_p$, choose
a distinct translation for each tuple of $\mathcal R$, and extend
all translations by fixed points on disjoint copies of
$I_1\sqcup\dotsb\sqcup I_r$.
Choose distinct core colors $c_1,\dotsc,c_r$. For the tuple
$(i_1,\dotsc,i_r)$, swap color $c_j$ with the entry at position
$i_j$, for each $j$.

Initially, every pair of words differs by one $p$-cycle. Take two
current words $u,v$ and consider their swaps for color $c_j$.
The target positions are fixed in
both current words, since this pool has not previously been used.
If the targets agree, applying the same transposition of colors
to both words leaves their relative permutation unchanged. If they
differ, say they are $a,b$, the new relative permutation is
\[
 (u^{-1}(c_j)\ b\ a)\,u^{-1}v.
\]
This inserts the two fixed points into its existing cycle.
The position $u^{-1}(c_j)$ lies on the cycle: the original translations put
$c_j$ in distinct core positions, and earlier swaps moved other colors.
Thus every pair still differs by exactly one nontrivial cycle.
The resulting family is a permutation matroid, since its prescribed
exchange interchanges the two parents. Its selected columns record
exactly the tuples of $\mathcal R$.

One core color was not selected. It occurs at a distinct core position
in every word and is unaffected by the swaps. Write $\lambda_t$ for
these actual assignment coordinates, one per word $u_t$. They are
the simplex barycentric coordinates. A complete binary description,
with the last equation for every cell $(e,c)$, is
\[
 \begin{gathered}
 \sum_t\lambda_t=1,\qquad
 \sum_{s\ne t}\lambda_s\leq1,\\
 x_{e,c}+\sum_{\{t:(u_t)_e\ne c\}}\lambda_t=1.
 \end{gathered}
\]
The first two conditions give nonnegative barycentric coordinates;
the last equations express every cell as their corresponding sum.
If a cell itself is $\lambda_t$, that coordinate is absent from
the sum in its equation, so all coefficients remain zero or one.
\end{proof}

\section[threeColumnObstruction]{Three columns and unbounded coefficients}

For an integer $m\geq1$, take the relation
\[
 \mathcal R_m=\{(i,j,i+j):0\leq i,j\leq m\}.
\]
Use three coordinate parts $a_0,\dotsc,a_m$,
$b_0,\dotsc,b_m$, and $c_0,\dotsc,c_{2m}$. Let $P_m$ be the
convex hull of
\[
 e_{a_i}+e_{b_j}+e_{c_{i+j}},\qquad0\leq i,j\leq m.
\]
By universality, $P_m$ is a three-column projection of a
binary-described permutation-matroid simplex.

\begin{theorem}[Affine equations and coefficient bounds]
The affine hull of $P_m$ is given by the three part sums equal to one
and the equation
\[
 \sum_{i=0}^m i\,x_{a_i}
 +\sum_{j=0}^m j\,x_{b_j}
 -\sum_{k=0}^{2m}k\,x_{c_k}=0.
\]
In particular, $\dim P_m=4m-1$. Every complete description of $P_m$
by integer equations and inequalities needs a coefficient of absolute
value at least $m$. The affine hull itself has minimum integer
coefficient radius exactly $m$.
\end{theorem}

\begin{proof}
An affine equation with coefficient arrays $\alpha_i,\beta_j,\gamma_k$
is constant with value $\delta$ on the vertices exactly when
\[
 \alpha_i+\beta_j+\gamma_{i+j}=\delta.
\]
Comparing $(i,j)$ with $(i-1,j+1)$ shows that every successive
difference of $\alpha$ and $\beta$ is the same number $t$.
Consequently all equations have the form
\[
 \begin{aligned}
 \alpha_i&=\alpha+ti,\\
 \beta_j&=\beta+tj,\\
 \gamma_k&=\gamma-tk.
 \end{aligned}
\]
These four parameters give exactly the four stated independent equations.
There are $4m+3$ coordinates, giving dimension $4m-1$.

For integral coefficients, $t$ is an integer. If $t\ne0$, the
$\gamma$ coefficients have range at least $2m$, requiring an absolute
value at least $m$. Hence any integer equation of smaller radius
is a combination of the three part sums. At a relative interior
point, the tight inequalities of a complete description must generate
all affine equations. Such smaller normals therefore cannot describe
$P_m$, even if inequalities as well as equations are allowed.

Finally, the fourth equation can be written
\[
 \begin{aligned}
 &\sum_{i=0}^m i\,x_{a_i}+\sum_{j=0}^m j\,x_{b_j}\\
 &\qquad+\sum_{k=0}^{2m}(m-k)x_{c_k}=m,
 \end{aligned}
\]
with all coefficients in $[-m,m]$.
\end{proof}

\subsection[ternaryWitness]{A point satisfying every ternary inequality}

For $m=2$, let
\[
 \begin{aligned}
 z_A=z_B&=(1/3,1/3,1/3),\\
 z_C&=(1/6,2/9,1/3,2/9,1/18).
 \end{aligned}
\]

\begin{proposition}[A ternary obstruction]
The point $z$ lies outside $P_2$ but satisfies every valid inequality
whose normal belongs to $\{-1,0,1\}^{11}$.
\end{proposition}

\begin{proof}
The average $\bar z$ of the nine vertices has
\[
 \begin{gathered}
 \bar z_C=(1,2,3,2,1)/9,\\
 z=\bar z+\frac1{18}(e_{c_0}-e_{c_4}).
 \end{gathered}
\]
For a ternary normal $a$, write
$h(a)=\max_{v\in\operatorname{vert}(P_2)}a\cdot v$.
If $a$ is constant on the vertices, the affine-equation classification
forces $t=0$, so its coefficients are constant within each part.
The perturbation then has score zero.

Otherwise the vertex scores are integers and not all equal, giving
\[
 h(a)-a\cdot\bar z
 =\frac19\sum_v(h(a)-a\cdot v)\geq\frac19.
\]
The perturbation changes the score by at most $2/18$, so
$a\cdot z\leq h(a)$. But the required weighted-balance equation
evaluates to $2/9$ at $z$, proving $z\notin P_2$.
\end{proof}

This example has a lift with $22$ positions and $22$ singleton colors:
take $p=11$ in the construction and add the $3+3+5$ pool positions.
The full polytope is an eight-dimensional simplex with a binary
description; its three-column image has dimension seven.
The \href{polytope.htm#generalPolytopeDescription}{coefficient question
for full assignment polytopes} remains a separate problem.
